\section{Experimental design}\label{sec:experiments}

\subsection{Instances}\label{sec:data}
\paragraph{Synthetic families.} We adapt the generator of \citet{cornuejols1991}, also used for
the \texttt{facilities} family of \citet{gasse2019}, to single sourcing:
\begin{itemize}
\item demand $d_j\sim U(5,35)$;
\item raw capacities $U(10,160)$, rescaled to a total-capacity-to-demand ratio
$r\in\{1.5, 3.0\}$;
\item fixed cost $f_i = (U(0,90) + U(100,110)\sqrt{Q_i})$;
\item service cost $c_{ij} = 10\, \mathrm{dist}_{ij}\, d_j$.
\end{itemize}
Customer locations follow three spatial families: uniform, clustered (Gaussian mixture) and
corridor (points along a few segments). Instances infeasible under single sourcing are rejected
by a first-fit-decreasing test followed, if needed, by a feasibility MILP. Splits are disjoint in
generator seeds; the corridor family never appears in training or validation.

\paragraph{External and real instances.} The 71 instances of \citet{holmberg1999}, with
published optima, use non-Euclidean costs and 10--30 facilities by 50--200 customers. The case
study uses delivered orders from the Olist public dataset aggregated by three-digit postal-code
prefix (up to 850 regions and 150 candidate facilities), with freight from the minimum road
freight tariff of the Brazilian regulator (ANTT); capacities and fixed costs are declared
modeling assumptions.

\begin{table}[t]\centering\small
\caption{Data sets. All files are frozen with SHA-256 hashes before any run.}\label{tab:data}
\resizebox{\linewidth}{!}{%
\begin{tabular}{llrll}
\toprule
Set & Size ($|I|\times|J|$) & Instances & Reference & Use \\
\midrule
Train & $30\times150$ & 160 & SCIP 60 s + LNS 30 s pool & GNN and selector training \\
Validation & $30\times150$ & 32 & same & parameter choices \\
Test & $30\times150$ & 32 & same & closed test \\
Corridor & $30\times150$ & 16 & same & unseen spatial family \\
Scale & $50\times200$ & 16 & SCIP 120 s + LNS 60 s & larger size \\
Holmberg & $10$--$30\times50$--$200$ & 71 & published optimum & external anchor \\
Olist & $30\times150$ to $150\times850$ & 4 & SCIP 300 s + LNS 120 s & real data \\
\bottomrule
\end{tabular}}
\end{table}

\subsection{Methods compared}
\begin{itemize}
\item \textbf{SCIP} 10.0 on the full model, one thread, default settings.
\item \textbf{Warm SCIP}: the LNS below for 10\% of the budget, then SCIP from its solution.
\item \textbf{LNS}: a matheuristic that frees customers with similar cost profiles or the
customers of one or two facilities and re-optimizes them with SCIP, with adaptive neighborhood
size.
\item \textbf{Adaptive expansion} alone (Section~\ref{sec:expansion}) over the full budget, with
the GNN score and, as a learning-free control, with the LP score \eqref{eq:lpscore}.
\item \textbf{Kernel search} \citep{guastaroba2014} over facilities (Section~\ref{sec:kernel}),
the classical counterpart of the expansion, with the same budget and solver.
\item \textbf{CLNS} from the LP-rounded start with each selector (rotation, ALNS, learned, GNN).
\item \textbf{Hybrid} (Algorithm~\ref{alg:hybrid}) with the GNN-guided, rotation and learned
selectors; \textbf{LP hybrid}, the same algorithm with the LP score in place of the GNN and the
rotation selector, in which no component is learned.
\item \textbf{Memory} variants (suffix \emph{+mem}) of CLNS and of both hybrids, with the
subproblem memory of Section~\ref{sec:memory}.
\item \textbf{Naive decomposition}: clusters solved independently with full capacities, then
greedily repaired.
\end{itemize}

\subsection{Protocol}\label{sec:protocol}
An audit of an earlier version of our harness found that solver set-up time was counted twice in
incumbent trajectories, that the primal integral did not use the cumulative best incumbent across
stages, and that there was no global deadline. The corrected protocol, frozen before the runs
reported here, is:
\begin{itemize}
\item \textbf{Equal wall-clock budget} per method (60~s for $30\times150$ and corridor, 120~s for
$50\times200$, 30~s for Holmberg), including feature computation, inference, LP, repair, model
building, solving and validation, under one monotonic clock.
\item \textbf{Isolation}: each run on a dedicated physical core (CPU affinity), with one core left
for the operating system, all numerical libraries limited to one thread, and imports and models
warmed up outside the timed region.
\item \textbf{Two clocks}: wall-clock and process CPU time. The ratio CPU/wall flags runs that
waited for the CPU; runs below 0.9 are marked as suspected contention.
\item \textbf{Frozen data}: instances and models are hashed in a versioned lock file; every run
verifies the lock and records the aggregate hash in its manifest.
\item \textbf{Seeds}: three solver seeds for stochastic methods, one for deterministic ones; seeds
are averaged within an instance before any statistics across instances. Three GNN checkpoints
were trained, but all optimization runs use the checkpoint of training seed 0; conclusions about
the learned ranking are limited to that checkpoint. Three GNN checkpoints
were trained, but all optimization runs use the checkpoint of training seed 0; conclusions about
the learned ranking are limited to that checkpoint.
\item \textbf{Independent validation}: every returned solution is re-evaluated from the original
data (assignment, capacity, objective).
\end{itemize}

\subsection{Metrics and statistics}\label{sec:metrics}
Let $z_k(t)$ be the cost of the best solution that run $k$ has found by wall-clock time $t$
(the cumulative minimum over all stages of the method; $z_k(t)=\infty$ before the first
solution), and $z_{\mathrm{BKS}}$ the best known cost of the instance: the minimum over all runs
of all methods and the reference label of Table~\ref{tab:data} (the published optimum for
Holmberg). We assume $z_{\mathrm{BKS}}>0$. Following the idea of the primal integral of
\citet{berthold2013}, with a gap function relative to the reference rather than to the larger of
the two objectives, we define the \emph{primal gap} and the normalized \emph{primal integral} as
\begin{equation}\label{eq:pi}
g_k(t)=\min\Bigl\{1,\ \frac{z_k(t)-z_{\mathrm{BKS}}}{z_{\mathrm{BKS}}}\Bigr\},
\qquad
P_k=\frac{1}{T}\int_0^T g_k(t)\,dt\ \in[0,1].
\end{equation}
$g_k$ is a non-increasing step function, so the integral is an exact sum over the incumbent
updates. $P_k=0$ means that the best known solution was available from time zero. A run with no solution
within the budget has $P_k=1$; the converse does not hold, since a solution more than 100\% above
the reference also has gap 1. The reference is computed within each run set and is not an
optimum; a method that finds its first solution, with a
1\% gap, at $t=6$~s and stays there has $P = 0.1\times 1+0.9\times0.01=0.109$ for $T=60$~s. The
primal integral is the primary metric because it rewards finding good solutions \emph{early},
which is the purpose of restricting the search space; the final gap alone does not distinguish a
method that is good after 5~s from one that is good after 55~s.

Secondary metrics are the \emph{final deviation} $g_k(T)$, computed from the best incumbent
with time stamp at most $T$ (a run without a solution by $T$ has deviation 1 and stays in the
denominator); the shares of instances with
$g(T)\le1\%$ and $g(T)\le0.1\%$; the time to reach $1\%$ (censored at $T$); for CLNS variants,
the \emph{repetition rate}, i.e., the fraction of selected subproblems whose key \eqref{eq:key}
had already failed with at least the same time limit; and the CPU/wall ratio of the run.

\paragraph{Aggregation and tests.} Seeds are averaged within an instance \emph{before} any
statistic across instances, so that the unit of analysis is the instance and stochastic methods
do not receive more weight than deterministic ones. Reported means carry 95\% bootstrap
confidence intervals over instances (4{,}000 resamples). Methods are compared with the two-sided paired
Wilcoxon signed-rank test on per-instance differences (zero differences discarded), reported
together with the mean difference and its bootstrap interval, with Holm's correction over the family of
comparisons of each hypothesis; we also report wins, ties and losses with a tolerance of
$10^{-4}$ on the deviation. Hypotheses, metrics and the choice $\alpha=0.5$ were pre-registered
before the corresponding runs.
